\FigureLayoutDeclare{img/2022/national_paper_a_science/q4_fig1.png}{7.0cm}{0bp 11bp 46bp 16bp}

\chapter{2022年全国甲卷（理）}

\section{单选题}

\begin{problem}
设 \(z=-1+\sqrt{3}\i\)，则 \(\frac{z}{z\bar z-1}=\)（\quad）

\choices
    {\(-1+3\i\)}
    {\(-1-3\i\)}
    {\(-\frac13+\frac{\sqrt3}{3}\i\)}
    {\(-\frac13-\frac{\sqrt3}{3}\i\)}
\end{problem}

\begin{answer}
C
\end{answer}

\begin{problem}
某社区通过公益讲座以普及社区居民的垃圾分类知识. 为了解讲座效果，随机抽取 \(10\) 位社区居民，让他们在讲座前和讲座后各回答一份垃圾分类知识问卷，这 \(10\) 位社区居民在讲座前和讲座后问卷答题的正确率如下图：

\bitmapfigure[max width=6.0cm]{img/2022/national_paper_a_science/q2_fig1.png}

则（\quad）

\choices
    {讲座前问卷答题的正确率的中位数小于 \(70\%\)}
    {讲座后问卷答题的正确率的平均数大于 \(85\%\)}
    {讲座前问卷答题的标准差小于讲座后正确率的标准差}
    {讲座后问卷答题的正确率的极差大于讲座前正确率的极差}
\end{problem}

\begin{answer}
B
\end{answer}

\begin{problem}
设全集 \(U=\{-2,-1,0,1,2,3\}\)，集合 \(A=\{-1,2\}\)，\(B=\{x\mid x^2-4x+3=0\}\)，则 \(\complement_U(A\cup B)\) 为（\quad）

\choices
    {\(\{1,3\}\)}
    {\(\{0,3\}\)}
    {\(\{-2,1\}\)}
    {\(\{-2,0\}\)}
\end{problem}

\begin{answer}
D
\end{answer}

\begin{problem}
如图，网格纸上绘制的是一个多面体的三视图，网格小正方形的边长为 \(1\)，则该多面体的体积为（\quad）

\bitmapfigure[width=7.0cm]{img/2022/national_paper_a_science/q4_fig1.png}

\choices
    {\(8\)}
    {\(12\)}
    {\(16\)}
    {\(20\)}
\end{problem}

\begin{answer}
B
\end{answer}

\begin{problem}
函数 \(y=\left( 3^x-3^{-x} \right)\cos x\) 在区间 \(\left[ -\frac{\pi}{2},\frac{\pi}{2} \right]\) 的图象大致为（\quad）

\choices
    {\choicebitmap[width=0.15\paperwidth]{img/2022/national_paper_a_science/q5_optA.png}}
    {\choicebitmap[width=0.15\paperwidth]{img/2022/national_paper_a_science/q5_optB.png}}
    {\choicebitmap[width=0.15\paperwidth]{img/2022/national_paper_a_science/q5_optC.png}}
    {\choicebitmap[width=0.15\paperwidth]{img/2022/national_paper_a_science/q5_optD.png}}
\end{problem}

\begin{answer}
A
\end{answer}

\begin{problem}
当 \(x=1\) 时，函数 \(f(x)=a\ln x+\frac b x\) 取得最大值 \(-2\)，则 \(f'(2)=\)（\quad）

\choices
    {\(-1\)}
    {\(-\frac12\)}
    {\(\frac12\)}
    {\(1\)}
\end{problem}

\begin{answer}
B
\end{answer}

\begin{problem}
在长方体 \(ABCD-A_1B_1C_1D_1\) 中，已知 \(B_1D\) 与平面 \(ABCD\) 和平面 \(AA_1B_1B\) 所成角均为 \(30^\circ\)，则（\quad）

\choices
    {\(AB=2AD\)}
    {\(AB\) 与平面 \(AB_1C_1D\) 所成的角为 \(30^\circ\)}
    {\(AC=CB_1\)}
    {\(B_1D\) 与平面 \(BB_1C_1C\) 所成的角为 \(45^\circ\)}
\end{problem}

\begin{answer}
D
\end{answer}

\begin{problem}
沈括的《梦溪笔谈》是中国古代科技史上的杰作，其中收录了计算圆弧长度的“会圆术”，如图，\(\myarc{AB}\) 是以 \(O\) 为圆心，\(OA\) 为半径的圆弧，\(C\) 是 \(AB\) 的中点，\(D\) 在 \(\myarc{AB}\) 上，\(CD\perp AB\)，“会圆术”给出的弧长的近似值 \(s\) 的计算公式：\(s=\overline{AB}+\frac{CD^2}{OA}\). 当 \(OA=2\)，\(\angle AOB=60^\circ\) 时，\(s=\)（\quad）

\bitmapfigure[max width=5cm]{img/2022/national_paper_a_science/q8_fig1.png}

\choices
    {\(\dfrac{11-3\sqrt3}{2}\)}
    {\(\dfrac{11-4\sqrt3}{2}\)}
    {\(\dfrac{9-3\sqrt3}{2}\)}
    {\(\dfrac{9-4\sqrt3}{2}\)}
\end{problem}

\begin{answer}
B
\end{answer}

\begin{problem}
甲、乙两个圆锥的母线长相等，侧面展开图的圆心角之和为 \(2\pi\)，侧面积分别为 \(S_{\text{甲}}\) 和 \(S_{\text{乙}}\)，体积分别为 \(V_{\text{甲}}\) 和 \(V_{\text{乙}}\). 若 \(\frac{S_{\text{甲}}}{S_{\text{乙}}}=2\)，则 \(\frac{V_{\text{甲}}}{V_{\text{乙}}}=\)（\quad）

\choices
    {\(\sqrt5\)}
    {\(2\sqrt2\)}
    {\(\sqrt{10}\)}
    {\(\dfrac{5\sqrt{10}}{4}\)}
\end{problem}

\begin{answer}
C
\end{answer}

\begin{problem}
椭圆 \(C:\frac{x^2}{a^2}+\frac{y^2}{b^2}=1\)（\(a>b>0\)）的左顶点为 \(A\)，点 \(P\)，\(Q\) 均在 \(C\) 上，且关于 \(y\) 轴对称. 若直线 \(AP\)，\(AQ\) 的斜率之积为 \(\frac14\)，则 \(C\) 的离心率为（\quad）

\choices
    {\( \dfrac{\sqrt3}{2}\)}
    {\( \dfrac{\sqrt2}{2}\)}
    {\( \dfrac{1}{2}\)}
    {\(\dfrac{1}{3}\)}
\end{problem}

\begin{answer}
A
\end{answer}

\begin{problem}
设函数 \(f(x)=\sin\left(\omega x+\frac\pi3\right)\) 在区间 \((0,\pi)\) 上恰有三个极值点，两个零点，则实数 \(\omega\) 的取值范围是（\quad）

\choices
    {\(\left[\frac53,\frac{13}6\right)\)}
    {\(\left[\frac53,\frac{19}6\right)\)}
    {\(\left(\frac{13}6,\frac83\right]\)}
    {\(\left(\frac{13}6,\frac{19}6\right]\)}
\end{problem}

\begin{answer}
C
\end{answer}

\begin{problem}
已知 \(a=\frac{31}{32}\)，\(b=\cos\frac14\)，\(c=4\sin\frac14\)，则（\quad）

\choices
    {\(c>b>a\)}
    {\(b>a>c\)}
    {\(a>b>c\)}
    {\(a>c>b\)}
\end{problem}

\begin{answer}
A
\end{answer}

\section{填空题}

\begin{problem}
设向量 \(\bs a\)，\(\bs b\) 的夹角的余弦值为 \(\frac13\)，且 \(|\bs a|=1\)，\(|\bs b|=3\)，则 \((2\bs a+\bs b)\cdot\bs b=\)\fillinblank{}
\end{problem}

\begin{answer}
11
\end{answer}

\begin{problem}
若双曲线 \(y^2-\frac{x^2}{m^2}=1\)（\(m>0\)）的渐近线与圆 \(x^2+y^2-4y+3=0\) 相切，则 \(m=\)\fillinblank{}
\end{problem}

\begin{answer}
\(\dfrac{\sqrt3}{3}\)
\end{answer}

\begin{problem}
从正方体的 \(8\) 个顶点中任选 \(4\) 个，则这 \(4\) 个点在同一个平面的概率为\fillinblank{}.
\end{problem}

\begin{answer}
\(\dfrac{6}{35}\)
\end{answer}

\begin{problem}
已知 \(\triangle ABC\)，点 \(D\) 在边 \(BC\) 上，\(\angle ADB=120^\circ\)，\(AD=2\)，\(CD=2BD\). 当 \(\frac{AC}{AB}\) 取得最小值时，\(BD=\)\fillinblank{}.
\end{problem}

\begin{answer}
\(\sqrt3-1\)
\end{answer}

\section{解答题}

\begin{problem}
记 \(S_n\) 为数列 \(\{a_n\}\) 的前 \(n\) 项和. 已知 \(\frac{2S_n}{n}+n=2a_n+1\).

\begin{enumerate}
\item 证明：\(\{a_n\}\) 是等差数列；
\item 若 \(a_4\)，\(a_7\)，\(a_9\) 成等比数列，求 \(S_n\) 的最小值.
\end{enumerate}
\end{problem}

\begin{solution}
\begin{enumerate}
\item 由已知等式可得
\[
S_n=na_n-\frac{n(n-1)}{2}
\]
当 \(n\ge2\) 时，分别对 \(n\) 和 \(n-1\) 使用上式，并利用 \(S_n-S_{n-1}=a_n\)，得
\[
a_n=na_n-(n-1)a_{n-1}-(n-1)
\]
所以
\[
a_n-a_{n-1}=1
\]
因此，\(\{a_n\}\) 是首项为 \(a_1\)、公差为 \(1\) 的等差数列.
\item 设 \(a_1=u\)，则 \(a_n=u+n-1\). 由 \(a_4\)，\(a_7\)，\(a_9\) 成等比数列，得
\[
(u+6)^2=(u+3)(u+8)
\]
解得 \(u=-12\)，从而 \(a_n=n-13\). 因此
\[
S_n=\frac{n(n+1)}{2}-13n=\frac{n(n-25)}{2}
\]
由于 \(n\) 为正整数，
\[
S_n+78=\frac{(n-12)(n-13)}{2}\ge0
\]
当 \(n=12\) 或 \(n=13\) 时等号成立，所以 \(S_n\) 的最小值为 \(-78\).
\end{enumerate}
\end{solution}

\begin{problem}
在四棱锥 \(P-ABCD\) 中，\(PD\perp\) 底面 \(ABCD\)，\(CD\parallel AB\)，\(AD=DC=CB=1\)，\(AB=2\)，\(DP=\sqrt3\).

\begin{enumerate}
\item 证明：\(BD\perp PA\)；
\item 求 \(PD\) 与平面 \(PAB\) 所成的角的正弦值.
\end{enumerate}

\bitmapfigure[max width=6.0cm]{img/2022/national_paper_a_science/q18_fig1.png}
\end{problem}

\begin{solution}
\begin{enumerate}
\item 在底面内建立直角坐标系，取 \(A=(0,0,0)\)，\(B=(2,0,0)\). 设 \(D=(u,v,0)\)，由 \(CD\parallel AB\) 且 \(CD=1\)，可设 \(C=(u+1,v,0)\). 由 \(AD=CB=1\)，有
\[
u^2+v^2=1
\]
和
\[
(u-1)^2+v^2=1
\]
两式相减得 \(u=\frac12\)，再由 \(v>0\) 得 \(v=\frac{\sqrt3}{2}\). 因为 \(PD\perp\) 平面 \(ABCD\) 且 \(DP=\sqrt3\)，所以
\[
D=\left(\frac12,\frac{\sqrt3}{2},0\right)
\]
以及
\[
P=\left(\frac12,\frac{\sqrt3}{2},\sqrt3\right)
\]
于是
\[
\overrightarrow{BD}=\left(-\frac32,\frac{\sqrt3}{2},0\right)
\]
和
\[
\overrightarrow{AP}=\left(\frac12,\frac{\sqrt3}{2},\sqrt3\right)
\]
因此
\[
\overrightarrow{BD}\cdot\overrightarrow{AP}=-\frac34+\frac34=0
\]
故 \(BD\perp PA\).
\item 设 \(PD\) 与平面 \(PAB\) 所成的角为 \(\theta\). 平面 \(PAB\) 的一个法向量为
\[
\overrightarrow{AB}\times\overrightarrow{AP}=(0,-2\sqrt3,\sqrt3)
\]
其长度为 \(\sqrt{15}\). 又 \(\lvert\overrightarrow{PD}\rvert=\sqrt3\)，所以
\[
\sin\theta
=\frac{\left|\overrightarrow{PD}\cdot\left(\overrightarrow{AB}\times\overrightarrow{AP}\right)\right|}{\lvert\overrightarrow{PD}\rvert\left\lvert\overrightarrow{AB}\times\overrightarrow{AP}\right\rvert}
=\frac{3}{\sqrt3\sqrt{15}}
=\frac{\sqrt5}{5}
\]
\end{enumerate}
\end{solution}

\begin{problem}
甲、乙两个学校进行体育比赛，比赛共设三个项目，每个项目胜方得 \(10\) 分，负方得 \(0\) 分，没有平局. 三个项目比赛结束后，总得分高的学校获得冠军. 已知甲学校在三个项目中获胜的概率分别为 \(0.5\)，\(0.4\)，\(0.8\)，各项目的比赛结果相互独立.

\begin{enumerate}
\item 求甲学校获得冠军的概率；
\item 用 \(X\) 表示乙学校的总得分，求 \(X\) 的分布列与期望.
\end{enumerate}
\end{problem}

\begin{solution}
\begin{enumerate}
\item 甲学校至少赢得两个项目时获得冠军，因此
\[
P(\text{甲赢两项})=0.5\times0.4\times0.2+0.5\times0.6\times0.8+0.5\times0.4\times0.8=0.44
\]
甲学校赢得三个项目的概率为
\[
P(\text{甲赢三项})=0.5\times0.4\times0.8=0.16
\]
所以
\[
P(\text{甲获冠军})=0.44+0.16=0.6
\]
\item 乙学校在三个项目中获胜的概率分别为 \(0.5\)，\(0.6\)，\(0.2\). 由各项目相互独立，得到 \(X\) 的分布列
\begin{center}
\begin{tabular}{c|cccc}
\(X\) & \(0\) & \(10\) & \(20\) & \(30\) \\ \hline
\(P\) & \(\dfrac4{25}\) & \(\dfrac{11}{25}\) & \(\dfrac{17}{50}\) & \(\dfrac3{50}\)
\end{tabular}
\end{center}
具体地，
\[
P(X=0)=0.5\times0.4\times0.8=\frac4{25}
\]
\[
P(X=10)=0.5\times0.4\times0.8+0.5\times0.6\times0.8+0.5\times0.4\times0.2=\frac{11}{25}
\]
\[
P(X=20)=0.5\times0.6\times0.8+0.5\times0.4\times0.2+0.5\times0.6\times0.2=\frac{17}{50}
\]
\[
P(X=30)=0.5\times0.6\times0.2=\frac3{50}
\]
所以
\[
E(X)=0\times\frac4{25}+10\times\frac{11}{25}+20\times\frac{17}{50}+30\times\frac3{50}=13
\]
\end{enumerate}
\end{solution}

\begin{problem}
设抛物线 \(C:y^2=2px\left( p>0 \right)\) 的焦点为 \(F\)，点 \(D\left( p,0 \right)\). 过 \(F\) 的直线交 \(C\) 于 \(M\)，\(N\) 两点. 当直线 \(MD\) 垂直于 \(x\) 轴时，\(|MF|=3\).

\begin{enumerate}
\item 求 \(C\) 的方程；
\item 设直线 \(MD\)，\(ND\) 与 \(C\) 的另一个交点分别为 \(A\)，\(B\)，记直线 \(MN\)，\(AB\) 的倾斜角分别为 \(\alpha\)，\(\beta\). 当 \(\alpha-\beta\) 取得最大值时，求直线 \(AB\) 的方程.
\end{enumerate}
\end{problem}

\begin{solution}
\begin{enumerate}
\item 抛物线 \(C\) 的焦点为 \(F\left(\frac p2,0\right)\). 因为 \(MD\perp x\) 轴，所以 \(M\) 的横坐标为 \(p\). 代入 \(y^2=2px\)，得 \(M\) 的纵坐标为 \(\pm p\sqrt2\). 因此
\[
|MF|=\sqrt{\left(p-\frac p2\right)^2+(\pm p\sqrt2)^2}=\frac{3p}{2}
\]
由 \(|MF|=3\) 得 \(p=2\)，故 \(C\) 的方程为 \(y^2=4x\).
\item 将 \(C\) 上的点参数化为 \(T(t)=(t^2,2t)\). 任取参数为 \(u\)，\(v\) 的两点，连接它们的直线方程为
\[
(u+v)y=2x+2uv
\]
将 \(T(u)\) 和 \(T(v)\) 的坐标分别代入上式可验证该直线方程.
由 \(M\)，\(N\) 所在直线过焦点 \(F=(1,0)\)，得 \(t\cdot t_N=-1\). 设 \(M=T(t)\)，则 \(N=T\left(-\frac1t\right)\).

直线 \(MD\) 过 \(D=(2,0)\)，由上面的弦方程得其与抛物线的两个交点参数之积为 \(-2\). 因而 \(A=T\left(-\frac2t\right)\). 同理，\(B=T(2t)\).

于是
\[
\tan\alpha=\frac{2}{t-\frac1t}=\frac{2t}{t^2-1}
\]
以及
\[
\tan\beta=\frac{2}{2t-\frac2t}=\frac{t}{t^2-1}
\]
令 \(r=\dfrac{t}{t^2-1}\). 当 \(r>0\) 时，
\[
\alpha-\beta=\arctan(2r)-\arctan r
\]
当 \(r<0\) 时，\(\alpha-\beta<0\)，而 \(t=\pm1\) 时两条直线均为竖直线，差值为 \(0\)，故最大值在 \(r>0\) 时取得. 对 \(r>0\)，令
\[
\varphi(r)=\arctan(2r)-\arctan r
\]
则
\[
\varphi'(r)=\frac2{1+4r^2}-\frac1{1+r^2}=\frac{1-2r^2}{(1+4r^2)(1+r^2)}
\]
所以 \(\varphi(r)\) 在 \(r=\frac1{\sqrt2}\) 时取得最大值. 此时
\[
t-\frac1t=\frac1r=\sqrt2
\]
直线 \(AB\) 对应的两个参数为 \(-\frac2t\)，\(2t\)，其和为 \(2\sqrt2\)，积为 \(-4\). 再由弦方程，得
\[
2\sqrt2y=2x-8
\]
即直线 \(AB\) 的方程为 \(x=\sqrt2y+4\).
\end{enumerate}
\end{solution}

\begin{problem}
已知函数 \(f(x)=\frac{\e^x}{x}-\ln x+x-a\).

\begin{enumerate}
\item 若 \(f(x)\ge 0\)，求 \(a\) 的取值范围；
\item 证明：若 \(f(x)\) 有两个零点 \(x_1\)，\(x_2\)，则 \(x_1x_2<1\).
\end{enumerate}
\end{problem}

\begin{solution}
\begin{enumerate}
\item 令
\[
g(x)=\frac{\e^x}{x}-\ln x+x
\]
其中 \(x>0\). 则
\[
g'(x)=\frac{\e^x(x-1)}{x^2}-\frac1x+1=\frac{(x-1)(\e^x+x)}{x^2}
\]
因为 \(x>0\) 时 \(\e^x+x>0\)，所以 \(g(x)\) 在 \((0,1)\) 上单调递减，在 \((1,+\infty)\) 上单调递增，最小值为
\[
g(1)=\e+1
\]
由 \(f(x)=g(x)-a\)，要使定义域内一切 \(x>0\) 都有 \(f(x)\ge0\)，充要条件是 \(a\le\e+1\). 因此 \(a\) 的取值范围为 \((-\infty,\e+1]\).
\item 若 \(f(x)\) 有两个零点，利用 \(g\) 在 \(1\) 的两侧分别严格单调，可设 \(0<x_1<1<x_2\)，且 \(g(x_1)=g(x_2)\). 令 \(x_1=\e^{-t}\)，其中 \(t>0\). 注意到
\[
g(\e^t)-g(\e^{-t})=\e^{\e^t-t}-\e^{\e^{-t}+t}+\e^t-\e^{-t}-2t
\]
设 \(h(t)=\e^t-\e^{-t}-2t\)，则 \(h(0)=0\)，且
\(h'(t)=\e^t+\e^{-t}-2>0\)，\((t>0)\).
故 \(\e^t-\e^{-t}>2t\). 从而
\[
\e^t-t>\e^{-t}+t
\]
这使得上式右端的前两项之差为正，后两项之和也为正，于是
\[
g\left(\frac1{x_1}\right)=g(\e^t)>g(\e^{-t})=g(x_1)=g(x_2)
\]
由于 \(g\) 在 \((1,+\infty)\) 上严格递增，得 \(x_2<\dfrac1{x_1}\)，所以 \(x_1x_2<1\).
\end{enumerate}
\end{solution}

\section{选考题：解答题}

\begin{problem}
在直角坐标系 \(xOy\) 中，曲线 \(C_1\) 的参数方程为 \(\begin{cases}
x=\dfrac{2+t}{6}\\
y=\sqrt t
\end{cases}\)（\(t\) 为参数），曲线 \(C_2\) 的参数方程为 \(\begin{cases}
x=-\dfrac{2+s}{6}\\
y=-\sqrt s
\end{cases}\)（\(s\) 为参数）.

\begin{enumerate}
\item 写出 \(C_1\) 的普通方程；
\item 以坐标原点为极点，\(x\) 轴正半轴为极轴建立极坐标系，曲线 \(C_3\) 的极坐标方程为 \(2\cos\theta-\sin\theta=0\)，求 \(C_3\) 与 \(C_1\) 交点的直角坐标，及 \(C_3\) 与 \(C_2\) 交点的直角坐标.
\end{enumerate}
\end{problem}

\begin{solution}
\begin{enumerate}
\item 由 \(y=\sqrt t\) 得 \(t=y^2\)，且 \(y\ge0\). 代入 \(x=\dfrac{2+t}{6}\)，得
\[
C_1:y^2=6x-2
\]
且 \(y\ge0\).
\item 设极坐标中的极径为 \(\rho\). 由 \(x=\rho\cos\theta\)，\(y=\rho\sin\theta\)，曲线 \(C_3\) 的方程化为
\[
2x-y=0
\]
即 \(y=2x\).

与 \(C_1\) 联立，得
\[
4x^2=6x-2
\]
所以 \(2x^2-3x+1=0\)，解得 \(x=\frac12\) 或 \(x=1\). 结合 \(y=2x\)，交点为
\[
\left(\frac12,1\right)
\]
和
\[
(1,2)
\]

由 \(C_2\) 的参数方程同理得 \(y^2=-6x-2\)，且 \(y\le0\). 与 \(y=2x\) 联立，得
\[
4x^2=-6x-2
\]
所以 \(2x^2+3x+1=0\)，解得 \(x=-\frac12\) 或 \(x=-1\). 因此 \(C_3\) 与 \(C_2\) 的交点为
\[
\left(-\frac12,-1\right)
\]
和
\[
(-1,-2)
\]
\end{enumerate}
\end{solution}

\begin{problem}
已知 \(a\)，\(b\)，\(c\) 均为正数，且 \(a^2+b^2+4c^2=3\)，证明：

\begin{enumerate}
\item \(a+b+2c\le 3\)；
\item 若 \(b=2c\)，则 \(\frac1a+\frac1c\ge 3\).
\end{enumerate}
\end{problem}

\begin{solution}
\begin{enumerate}
\item 由柯西不等式，
\[
(a+b+2c)^2\le(1^2+1^2+1^2)(a^2+b^2+4c^2)=9
\]
因为 \(a\)，\(b\)，\(c\) 均为正数，所以 \(a+b+2c>0\)，从而 \(a+b+2c\le3\).
\item 当 \(b=2c\) 时，条件化为 \(a^2+8c^2=3\). 由柯西不等式，
\[
\left(\frac1a+\frac1c\right)(a+4c)\ge(1+2)^2=9
\]
另一方面，
\[
a+4c=a\cdot1+(\sqrt8c)\cdot\frac4{\sqrt8}\le\sqrt{(a^2+8c^2)\left(1+\frac{16}{8}\right)}=3
\]
由于 \(a+4c>0\)，因此
\[
\frac1a+\frac1c\ge\frac9{a+4c}\ge3
\]
\end{enumerate}
\end{solution}
